\section{Elementary integral chain constructions}
\label{mapsI:chain:section}

These classical constructions are included as rederivations of the singular-chain prerequisites; they make no novelty claim.

\subsection{Chains, pairs and homotopies}
\label{mapsI:chain:section-1}

Let \(C_m(X)\) be the free abelian group on the continuous maps \(\sigma:\Delta^m\to X\), with \(C_m=0\) for \(m<0\). Write \(\sigma[i_0,\ldots,i_k]\) for its restriction to the indicated ordered face. The differential is

\[
\partial\sigma=\sum_{i=0}^m(-1)^i\sigma[0,\ldots,\widehat i,\ldots,m].
\]

Each twice-deleted face appears with opposite signs in \(\partial^2\sigma\), so \(\partial^2=0\). Homology is \(H_m=\ker\partial/\operatorname{im}\partial\). For \(A\subset X\), the relative complex is \(C_*(X,A)=C_*(X)/C_*(A)\). It is free on the simplices not entirely contained in \(A\). With an abelian coefficient group \(G\), use \(C_*(X,A)\otimes G\) for homology and \(\operatorname{Hom}(C_*(X,A),G)\) for cohomology; its coboundary is \(\delta a=a\partial\). Maps of pairs induce the evident chain maps and contravariant cochain maps.

\begin{lemma}[homotopy invariance]
\label{mapsI:chain:lemma:1.1}
Homotopic maps of pairs induce the same homology and cohomology maps, with every coefficient group.
\end{lemma}

\begin{proof}
For a homotopy \(H:X\times I\to Y\), triangulate \(\Delta^m\times I\) into the \(m+1\) ordered simplices

\[
[(v_0,0),\ldots,(v_i,0),(v_i,1),\ldots,(v_m,1)],
\qquad 0\leq i\leq m.
\]

Apply \(H(\sigma(-),-)\) to their sum with coefficients \((-1)^i\), defining \(P\sigma\). Adjacent internal faces cancel. The bottom and top faces remain with signs minus and plus, and the side faces are \(-P\partial\sigma\). Hence
\(\partial P+P\partial=H_{1*}-H_{0*}\).
If \(H(A\times I)\subset B\), this formula descends to relative chains. Tensoring the homotopy, or precomposing cochains with it, proves the assertion.
\end{proof}

The point complex has one generator in every nonnegative degree; its differential is zero in odd degrees and the identity in positive even degrees. Therefore a point has homology \(\mathbb Z\) in degree zero and zero elsewhere. By \cref{mapsI:chain:lemma:1.1} the same is true of a contractible space. Moreover \(H_0(X)\) is free on the path components: the boundary of a path equates its two endpoints, and these are exactly the relations imposed by paths.

\subsection{Making singular chains small}
\label{mapsI:chain:section-2}

For an open cover \(\mathcal U\) of \(X\), let \(C_*^{\mathcal U}(X)\) be generated by simplices whose image lies in one cover member. Local arguments require the inclusion of these chains to be a chain homotopy equivalence, rather than just an intuitive subdivision picture.

\begin{lemma}[small-chain lemma]
\label{mapsI:chain:lemma:2.1}
There is a chain map \(\rho:C_*(X)\to C_*^{\mathcal U}(X)\) and a homotopy \(D\) with

\[
\rho\iota=1,\qquad 1-\iota\rho=\partial D+D\partial.
\]

Both preserve every subspace of \(X\), so the assertion also holds for relative complexes.
\end{lemma}

\begin{proof}
First construct subdivision \(S\) on the affine chains in each standard simplex. If \(b\) is its barycentre, define \(S\sigma=b*(S\partial\sigma)\) inductively, starting with vertices. Here the affine cone with apex first satisfies
\(\partial(b*c)=c-b*(\partial c)\) for a positive-dimensional cycle, and the corresponding reduced formula for dimension zero. It follows inductively that \(\partial S=S\partial\). Composing these affine subdivisions with a singular simplex defines a natural chain map on singular chains.

Construct a subdivision homotopy \(T\) by induction too. On an \(m\)-simplex, the chain
\(z=\sigma-S\sigma-T\partial\sigma\) is a cycle by the already established lower-dimensional identity. It has augmentation zero when relevant. Cone it to the barycentre and set \(T\sigma=b*z\). Then
\(\partial T+T\partial=1-S\).
All chains constructed remain inside the original simplex image. Thus subdivision and its homotopy preserve small chains and chains in any subspace.

A barycentric subsimplex has vertices that are barycentres of nested faces. If a face of size \(\ell\) lies in one of size \(k\), their barycentre difference has norm at most \((k-\ell)/k\) times the old diameter, at most \(m/(m+1)\) times it. The diameter of a Euclidean simplex is the maximum of its vertex distances. Iteration therefore makes the mesh tend to zero. For the inverse-image open cover of the compact standard simplex, choose at each point a ball whose doubled ball lies in one cover member. Finitely many of the smaller balls cover the simplex; the minimum of their radii is positive. Every subset of smaller diameter, by choosing one of its points in a smaller ball, lies in the corresponding doubled ball. This is the required Lebesgue number. Every singular simplex, and every finite chain, becomes small after sufficiently many subdivisions.

Put \(T_N=\sum_{j=0}^{N-1}S^jT\); then
\(\partial T_N+T_N\partial=1-S^N\).
Now construct \(\rho,D\) inductively on the basis of singular simplices. Given their values on the boundary of \(\sigma\), put
\(z=\sigma-D\partial\sigma\), so \(\partial z=\rho\partial\sigma\) is small. Choose \(N\) such that \(S^Nz\) is small and define

\[
\rho\sigma=S^Nz+T_N(\rho\partial\sigma),
\qquad D\sigma=T_Nz.
\]

The second term of \(\rho\sigma\) is small because \(T_N\) preserves smallness. The displayed identities give
\(\partial\rho\sigma=\rho\partial\sigma\) and
\(\sigma-\rho\sigma=\partial D\sigma+D\partial\sigma\).
If \(\sigma\) is already small, its faces have \(D=0\) by induction; choose \(N=0\), obtaining \(\rho\sigma=\sigma\), \(D\sigma=0\). Every operation remains in the original image, which proves the subspace assertion.
\end{proof}

\begin{corollary}[excision]
\label{mapsI:chain:corollary:2.2}
If \(Z\subset A\subset X\) and \(\overline Z\subset\operatorname{int}_X A\), the inclusion
\((X\setminus Z,A\setminus Z)\to(X,A)\)
induces homology and cohomology isomorphisms with every coefficient group.
\end{corollary}

\begin{proof}
Use the open cover \(\operatorname{int}A, X\setminus\overline Z\) of \(X\). After quotienting small chains by those in \(A\), the first summand disappears and the quotient is
\(C_*(X\setminus\overline Z,A\setminus\overline Z)\).
The restricted cover of \(X\setminus Z\) gives the very same quotient: its first member is \((\operatorname{int}A)\setminus Z\), whose chains again disappear. \cref{mapsI:chain:lemma:2.1} identifies both ordinary relative complexes with these quotients up to chain homotopy. Tensor and dualize the equivalences.
\end{proof}

\subsection{Exact sequences and local gluing}
\label{mapsI:chain:section-3}

A degreewise exact sequence of complexes \(0\to K\to L\to M\to0\) has a long exact sequence in homology. Its connecting map sends a cycle \(m\) to the class of \(\partial\ell\) in \(K\), where \(\ell\) lifts \(m\). Changing the lift changes this by a boundary; changing the cycle by a boundary does likewise. Exactness follows directly: a cycle of \(L\) maps to a boundary of \(M\) exactly when, after subtracting the boundary of a lift, it lies in \(K\); a cycle of \(M\) has zero connecting class exactly when its lift can be corrected by a chain in \(K\) to a cycle of \(L\); a cycle of \(K\) becomes a boundary in \(L\) exactly when it is the connecting image of the projected bounding chain. These three arguments prove exactness at all the repeating positions. The same lifting argument, with \(\delta\), proves the cohomological version and its naturality.

The pair sequence comes from
\(0\to C_*(A)\to C_*(X)\to C_*(X,A)\to0\).
It is degreewise split on the simplex bases, so it stays exact after tensoring or dualizing.

For an open cover \(X=U\cup V\), the sequence

\[
0\to C_*(U\cap V)\xrightarrow{c\mapsto(c,-c)}
C_*(U)\oplus C_*(V)\xrightarrow{(a,b)\mapsto a+b}
C_*^{\{U,V\}}(X)\to0
\]

is degreewise split exact: each simplex belongs to \(U\), to \(V\), or to both. \cref{mapsI:chain:lemma:2.1} replaces the last complex by \(C_*(X)\). We obtain the \textbf{Mayer--Vietoris sequences}. For cohomology the middle-to-intersection map is the difference of restrictions. Dividing throughout by the corresponding chains in a subspace \(A\) gives the relative sequences for \((X,A)\), with pieces \((U,U\cap A)\), \((V,V\cap A)\) and \((U\cap V,U\cap V\cap A)\). Splitting is again obtained on the simplices outside \(A\).

We will use the following precise form of the five lemma. In a commutative diagram of exact five-term sequences, if maps at positions one, two, four and five are isomorphisms, the middle map is an isomorphism. For injectivity, an element with zero image maps to zero at position four, so comes from position two; lift its image-zero condition through position one and use the isomorphisms there to correct its preimage to zero. For surjectivity, map a target middle element to position four; lift there, note its image at position five vanishes, and lift back to a source middle element. Their difference lies in the image of position two, whose isomorphism provides the required correction. This proves the lemma and applies equally to homology and cohomology segments.

For example, cover \(S^m\), \(m\geq1\), by slightly enlarged hemispheres. Each is contractible, and their intersection retracts to \(S^{m-1}\). The reduced Mayer--Vietoris sequence gives
\(\widetilde H_j(S^m;G)\cong\widetilde H_{j-1}(S^{m-1};G)\).
Starting with two points, this yields \(\widetilde H_m(S^m;G)=G\) and zero in other degrees; the same argument gives reduced cohomology. The generator obtained by orienting the two hemispheres compatibly is the fundamental class of the oriented sphere. Concretely triangulate it as the boundary of an \(m+1\)-simplex, radially projected from an interior point. Each ray meets the boundary once, and its intersection distance is continuous face by face, so this is a homeomorphism. The alternating sum of oriented faces is a cycle because shared faces cancel. Its image at a point inside one face is the positive local disk generator by excision. Since the sphere homology group in degree \(m\) is infinite cyclic, that value one proves the cycle is a generator. To verify the same local value at any point, choose a different centre or a rotated simplex whose face interior contains that point; the radial projections are related through orientation-preserving rotations, which are homotopic to the identity by the following elementary plane-rotation argument: an orthogonal orientation-preserving matrix is reduced to the identity by successive rotations in coordinate two-planes, each connected to the identity by its angle parameter. This gives the fundamental class \([S^m]\), with local value one in every oriented chart.

\subsection{Products, signs and evaluation on fibres}
\label{mapsI:chain:section-4}

For cochains of degrees \(p,q\) with values in a commutative ring \(R\), define the Alexander--Whitney cup product by

\[
(a\smile b)(\sigma)=a(\sigma[0,\ldots,p])
b(\sigma[p,\ldots,p+q]).
\]

Expanding boundaries, terms deleting vertices strictly before the joining vertex give \(\delta a\smile b\), those after it give \((-1)^p a\smile\delta b\), and the two joining terms cancel. Thus

\[
\delta(a\smile b)=\delta a\smile b+(-1)^p a\smile\delta b.
\]

The formula is natural. It is associative: the same three face intervals occur in either order of multiplication. The constant zero-cochain \(1\) is its unit. It descends to cohomology. An absolute class can multiply a relative class for every pair; cochains on the subspace vanish in the relative factor.

We supply the chain comparison that explains external products and graded commutativity. On the total tensor complex use
\(\partial(c\otimes d)=\partial c\otimes d+(-1)^{\deg c}c\otimes\partial d\).
Define

\[
A\sigma=\sum_{p+q=m}\sigma_X[0,\ldots,p]\otimes\sigma_Y[p,\ldots,m],
\qquad \sigma:\Delta^m\to X\times Y.
\]

Conversely define \(S(c\otimes d)\) by triangulating \(\Delta^p\times\Delta^q\) along all lattice paths from \((0,0)\) to \((p,q)\). Give a path the sign of the permutation of its horizontal and vertical steps relative to all horizontal steps followed by all vertical steps. Compose its affine simplex with \((c,d)\). Internal faces from paths interchanging adjacent unlike steps cancel; boundary faces give precisely \(\partial c\otimes d+(-1)^p c\otimes\partial d\). For \(A\), deleting a nonjoining vertex gives the matching tensor boundary and joining-vertex terms telescope. Hence both maps are chain maps, preserving augmentation.

\begin{lemma}[product chain equivalence]
\label{mapsI:chain:lemma:4.1}
\leavevmode\par
The two maps are inverse up to natural chain homotopies, preserving products of subspaces.
\end{lemma}

\begin{proof}
Here is an explicit inductive construction of the needed homotopies. Chains in a convex set have a contraction to any fixed vertex: cone each singular simplex to that vertex by linear interpolation. The cone boundary identity gives \(\partial K+K\partial=1-i\epsilon\), where \(\epsilon\) is augmentation and \(i\) inserts the vertex. Tensoring two such contractions gives a contraction of their tensor complex to \(\mathbb Z\): the homotopy is \(K\otimes1+(i\epsilon)\otimes L\), with the tensor sign \((-1)^{\deg c}\) in the second term. The first identity and then the second verify this formula; its boundary sum is \(1-(i\epsilon)\otimes(j\epsilon)\).

For \(1-SA\), assume a homotopy \(h\) is already defined on lower-dimensional simplices. On the universal diagonal simplex \(\Delta^m\to\Delta^m\times\Delta^m\), form
\(z=(1-SA)\sigma-h\partial\sigma\).
It is an augmented cycle. Cone \(z\) in the convex product to fill it. For an arbitrary \(\sigma\), push this filling forward by \(\sigma_X\times\sigma_Y\). This defines \(h\) naturally and gives \(\partial h+h\partial=1-SA\).

For \(1-AS\), repeat on a universal tensor generator \(1_{\Delta^p}\otimes1_{\Delta^q}\). The residual cycle is filled by the tensor contraction just constructed, then pushed forward by the two singular simplices. Induction on \(p+q\) gives the other homotopy. If either simplex lies in a subspace, every filling still maps into the corresponding product subspace.
\end{proof}

The same induction compares \(A\) with the map obtained by swapping the two coordinates and then applying \(A\) and the signed tensor swap
\(c_p\otimes d_q\mapsto(-1)^{pq}d_q\otimes c_p\).
They agree on augmentation, and the tensor contraction fills the residual cycles exactly as in \cref{mapsI:chain:lemma:4.1}. Evaluate \(a\otimes b\) on this homotopy after the diagonal \(X\to X\times X\). For cocycles the resulting difference is a coboundary. Therefore

\[
[a]\smile[b]=(-1)^{pq}[b]\smile[a].
\]

No assertion of cochain commutativity was used.

The equivalence also descends to
\(C_*(X,A)\otimes C_*(Y,B)\) and the product complex modulo
\(C_*(A\times Y)+C_*(X\times B)\).
When \(A,B\) are open, small chains in their union identify this quotient up to chain homotopy with
\(C_*(X\times Y,A\times Y\cup X\times B)\).
If one subspace is empty the assertion holds for the other subspace without openness. For classes, define the external product by \(a\times b=\operatorname{pr}_1^*a\smile\operatorname{pr}_2^*b\); its cochain formula is evaluation of \(a\otimes b\) on \(A\). This defines external products on all pairs used here. Cup products of two relative classes likewise take values relative to their union when its two subspaces are open, by the same small-chain comparison.

For an integral relative cocycle \(u\) of degree \(r\), define a right cap operator on an \(m\)-simplex, \(m\geq r\), by

\[
R_u\sigma=u(\sigma[m-r,\ldots,m])\,\sigma[0,\ldots,m-r],
\]

and set it to zero when \(m<r\). Boundary expansion and \(\delta u=0\) give
\(\partial R_u=R_u\partial\); the faces at the cut cancel by the cocycle equation on the last \(r+1\) vertices. Thus it is a degree-\(-r\) map commuting with the unshifted differentials. If \(u\) vanishes on a subspace, this map descends from relative chains to absolute chains. For a cocycle \(a\) of degree \(m-r\), its defining identity is

\[
a(R_u\sigma)=(a\smile u)(\sigma).
\]

We use this convention throughout. Reindexing chain groups here keeps their differential unchanged; it is not the alternative signed convention for suspended chain complexes. Replacing \(u\) by \(u+\delta v\) changes the cap map by a homotopy: on \(C_m\) set \(h=(-1)^{m-r+1}R_v\). The cup-product differential identity gives \(\partial h+h\partial=R_{\delta v}\). Products and this cap identity also make sense with an abelian coefficient group in the first factor and an integral cochain in the second, using integer multiplication on the group.
